\chapter{Números, fechas y horas}
\section{Objetivos}
\begin{itemize}[leftmargin=*]
	\item Contar del 0 al 100 y manejar precios simples.
	\item Decir la fecha, el día y la hora en formato neerlandés.
	\item Preguntar y dar la hora con precisión de minutos.
\end{itemize}

\section{Contenidos clave}
\begin{itemize}[leftmargin=*]
	\item Números cardinales 0--100, decenas y centenas básicas.
	\item Meses, días de la semana, estaciones; preposición \textit{in/op} para fechas.
	\item Hora: formato de media hora/\textit{over/voor half}, preguntar \textit{Hoe laat is het?}.
	\item Precios con euro y céntimos.
\end{itemize}

\section{Ruta del capítulo}
Siete secciones cortas con listas, patrones y práctica.

\section{Números 0\textendash{}100}
\textbf{Patrón base.} Del 0 al 20 se memorizan; a partir del 21 se construyen ``unidad + en + decena''.

\begin{tabular}{@{}lll@{}}
\toprule
0--10 & \textit{nul, een, twee, drie, vier, vijf, zes, zeven, acht, negen, tien} &  \\
11--20 & \textit{elf, twaalf, dertien, veertien, vijftien, zestien, zeventien, achttien, negentien, twintig} &  \\
Tens & \textit{dertig, veertig, vijftig, zestig, zeventig, tachtig, negentig, honderd} &  \\
\bottomrule
\end{tabular}

\textbf{Construcción 21\textendash{}99}
\begin{itemize}[leftmargin=*]
  \item Forma: unidad + \textit{en} + decena. Ej.: 24 = \textit{vierentwintig}, 57 = \textit{zevenenvijftig}.
  \item Excepción menor: 1 = \textit{een} mantiene la raíz \textit{een} sin acento.
\end{itemize}

\textbf{Mini práctica}
\begin{itemize}[leftmargin=*]
  \item Escribe y lee en voz alta: 13, 28, 46, 73, 99.
  \item Dictado inverso: tu compañero dice el número en neerlandés, escribes la cifra.
\end{itemize}

\section{Números grandes básicos}
\textbf{Centenas y miles.}
\begin{itemize}[leftmargin=*]
  \item 100 = \textit{honderd}; 200 = \textit{tweehonderd}; 999 = \textit{negenhonderdnegenennegentig}.
  \item 1.000 = \textit{duizend}; 2.500 = \textit{tweeduizend vijfhonderd}.
\end{itemize}

\textbf{Reglas rápidas}
\begin{itemize}[leftmargin=*]
  \item Usa espacio entre bloques grandes: \textit{tweeduizend vijfhonderd}. 
  \item Para años se suele omitir \textit{en}: 1987 = \textit{negentienhonderdzevenentachtig}. 
  \item Para precios altos: 1.299 = \textit{twaalfhonderd negenennegentig} también es aceptado.
\end{itemize}

\textbf{Mini práctica}
\begin{itemize}[leftmargin=*]
  \item Escribe 345, 1.050 y 2.021 en neerlandés.
  \item Lee en voz alta 7.800 y 15.000.
\end{itemize}

\section{Días de la semana}
\textbf{Lista.} \textit{maandag, dinsdag, woensdag, donderdag, vrijdag, zaterdag, zondag}.

\textbf{Uso}
\begin{itemize}[leftmargin=*]
  \item Preposición: \textit{op maandag} (el lunes), a menudo se omite en habla coloquial: \textit{maandag ga ik werken}.
  \item Plural para hábitos: \textit{Op zaterdagen sport ik} (menos frecuente; se puede usar singular).
\end{itemize}

\textbf{Mini práctica}
\begin{itemize}[leftmargin=*]
  \item Escribe tu agenda de la semana usando \textit{op + día}.
  \item Escucha y marca el día mencionado: el profe dice tres frases.
\end{itemize}

\section{Meses del año}
\textbf{Lista.} \textit{januari, februari, maart, april, mei, juni, juli, augustus, september, oktober, november, december}.

\textbf{Uso}
\begin{itemize}[leftmargin=*]
  \item Fechas con mes: \textit{in mei}, \textit{in december}.
  \item Con día y mes: \textit{op 3 mei}, \textit{op 21 juli}.
  \item Sin mayúscula inicial en neerlandés.
\end{itemize}

\textbf{Mini práctica}
\begin{itemize}[leftmargin=*]
  \item Escribe tu mes de nacimiento: \textit{Ik ben jarig in \\underline{\hspace{1.2cm}}}.
  \item Lee tres fechas con día y mes en voz alta.
\end{itemize}

\section{Fechas}
\textbf{Formato hablado.} Día + mes (sin artículo) y \textit{van} para decir la fecha completa: \textit{Het is vijf juni}. Año añadido: \textit{Het is vijf juni tweeduizendvierentwintig}.

\textbf{Orden y preposición}
\begin{itemize}[leftmargin=*]
  \item Para eventos: \textit{op 14 februari}, \textit{op 1 mei}. 
  \item Para meses/estaciones: \textit{in augustus}, \textit{in de winter}.
  \item Ordinales especiales: 1 = \textit{de eerste}, 2 = \textit{de tweede}, 3 = \textit{de derde}. 
\end{itemize}

\textbf{Mini práctica}
\begin{itemize}[leftmargin=*]
  \item Di en voz alta tres fechas: 1/5, 14/2, 25/12.
  \item Escribe cuándo es tu próxima cita: \textit{Mijn afspraak is op \\underline{\hspace{1.2cm}}}. 
\end{itemize}

\section{La hora (formal e informal)}
\textbf{Preguntar y responder.} \textit{Hoe laat is het? Het is drie uur}. Para minutos: \textit{Het is vijf over drie}.

\textbf{Patrones informales frecuentes}
\begin{itemize}[leftmargin=*]
  \item \textit{over} = minutos después: \textit{tien over twee} (2:10).
  \item \textit{voor} = minutos antes: \textit{vijf voor zes} (5:55).
  \item \textit{half} = media hora antes de la hora siguiente: \textit{half vier} = 3:30.
  \item Combinados: \textit{vijf over half vijf} (4:35), \textit{tien voor half negen} (8:20).
\end{itemize}

\textbf{Formalidad y 24h}
\begin{itemize}[leftmargin=*]
  \item Reloj de 24h en horarios: 18:45 = \textit{achttien uur vijfenveertig}.
  \item Conversación informal prefiere patrones con \textit{half/over/voor}.
\end{itemize}

\textbf{Mini práctica}
\begin{itemize}[leftmargin=*]
  \item Escribe y pronuncia: 7:15, 9:40, 12:05, 18:30.
  \item Convierte de formato 24h a informal: 14:20, 21:50.
\end{itemize}

\section{Horarios y citas}
\textbf{Hablar de agenda.} Combina día + hora + lugar: \textit{We hebben les op dinsdag om acht uur}. 

\textbf{Frases útiles}
\begin{itemize}[leftmargin=*]
  \item Proponer: \textit{Zullen we om vier uur afspreken?}
  \item Confirmar: \textit{Dat kan. Tot dan!}
  \item Reprogramar: \textit{Kan het om half negen?}
  \item Cancelar cortésmente: \textit{Sorry, ik kan niet komen. Zullen we morgen?}
\end{itemize}

\textbf{Mini práctica}
\begin{itemize}[leftmargin=*]
  \item Escribe tres citas reales de tu semana usando día y hora.
  \item Role-play: proponer, mover y cancelar una cita.
\end{itemize}


\section{Vocabulario clave}
\begin{itemize}[leftmargin=*]
	\item Días: \textit{maandag, dinsdag, woensdag, donderdag, vrijdag, zaterdag, zondag}.
	\item Meses: \textit{januari} \ldots{} \textit{december}.
	\item Hora: \textit{uur}, \textit{half}, \textit{kwart}, \textit{over}, \textit{voor}, \textit{middag}, \textit{avond}, \textit{'s ochtends}.
\end{itemize}

\section{Pronunciación}
\begin{itemize}[leftmargin=*]
	\item Alarga vocal en \textit{vier}, \textit{twee}; evita diptongar.
	\item Marca /x/ suave en \textit{acht}, \textit{acht-tien}; no añadas vocal extra.
	\item En hora informal, \textit{half drie} suena \textipa{/'hAlf dri/}; sin /e/ final.
\end{itemize}

\section{Errores comunes}
\begin{itemize}[leftmargin=*]
	\item Invertir decenas/unidades: \textit{vijf-en-twintig} (25), no \textit{twintig-en-vijf}.
	\item Decir artículo con fecha: usa \textit{op 3 mei}, sin artículo.
	\item Confundir \textit{half drie} (2:30) con 3:30.
\end{itemize}

\section{Práctica guiada}
\begin{itemize}[leftmargin=*]
	\item Diccionarios inversos: escuchas números y los escribes.
	\item Escribe 10 fechas personales y conviértelas a ordinal + mes en neerlandés.
	\item Empareja horarios escritos con relojes analógicos (informal vs. 24h).
\end{itemize}

\section{Ejercicios de producción}
\begin{itemize}[leftmargin=*]
	\item Graba un mensaje dejando tu número, cita y hora de encuentro.
	\item Escribe una agenda semanal breve con 8 actividades y horas.
\end{itemize}

\section{Mini repaso}
\begin{itemize}[leftmargin=*]
	\item Lees y dices números hasta 100 con buena pronunciación.
	\item Expresas fechas con \textit{op} y ordinal correcto.
	\item Das la hora en formato formal e informal sin confusiones.
\end{itemize}

\section{Resultado esperado}
\begin{itemize}[leftmargin=*]
	\item Al terminar puedes dar y entender números, fechas y horas cotidianas, agendar citas y verificar horarios sin ambigüedad.
\end{itemize}

\section{Práctica sugerida}
\begin{itemize}[leftmargin=*]
	\item Bingo de números y dictados de teléfonos/precios.
	\item Calendario: marcar eventos y leer fechas en voz alta.
	\item Reloj analógico: preguntas rápidas con minutos.
\end{itemize}

\section{Microevaluación}
\begin{itemize}[leftmargin=*]
	\item Audio de 10 cifras/horas: escribe lo que oyes y léelo después.
\end{itemize}
